\section{Methodology}
\label{sec:methodology}

\subsection{Problem Formulation}
\label{sec:formulation}

\paragraph{Sets and parameters.}
Let $\mathcal{S}$, $\mathcal{A}$, $\mathcal{D}$, $\mathcal{C}$, $\mathcal{I}$, $\mathcal{L}$, $\mathcal{E}$, $\mathcal{G}$ denote scenes, actors, days, cities, towns, locations, equipment types, and mandatory co-scheduling groups.
Surjective mappings $\phi: \mathcal{L} \to \mathcal{I}$ and $\psi: \mathcal{I} \to \mathcal{C}$ encode the three-tier geographic hierarchy.
Each scene $s \in \mathcal{S}$ has shooting duration $\delta_s$, required actors $A(s) \subseteq \mathcal{A}$, compatible locations $L(s) \subseteq \mathcal{L}$, and an optional time window with deadline penalty.
Each actor $a \in \mathcal{A}$ has daily wage $w_a$, availability calendar $\mathcal{D}_a$, consecutive-day limit $\bar{h}_a$, booking window $[\beta^s_a, \beta^e_a]$, and holding fee $h_a$.

\paragraph{Decision variables.}
Seven binary families: $x_{s,d}$ (scene--day assignment), $y_{s,l}$ (scene--location assignment), $v_{a,s,d}$ (actor participation), $z_{l,s,d}$ (location--scene--day), $u_{i,d}$ (active town per day), $\phi_{a,d}$ (actor idle on booked day), $w_{a,d}$ (actor working-day indicator).

\paragraph{Bi-objective formulation.}
We simultaneously minimise cost $f_1$ and makespan $f_2$:
\begin{align}
f_1 &= \underbrace{\sum_{a} w_a \bigl|\{d : \sum\nolimits_{s} v_{a,s,d} \geq 1\}\bigr|}_{\text{actor wages}}
     + \underbrace{\sum_{d}\;\sum_{l \neq l'} c_{l,l'}\,\sigma_{l,l',d}}_{\text{hierarchical transfers}}
     + \underbrace{\sum_{s} \rho_s\,(\mathrm{day}(s) - \bar{d}_s)^{+}}_{\text{deadline penalties}}
     \nonumber \\
    &\quad
     + \underbrace{\sum_{l,s,d} \frac{z_{l,s,d}}{\kappa^{L}_{l}}
                 + \sum_{a,s,d} \frac{v_{a,s,d}}{\kappa^{A}_{a}}}_{\text{criticality-weighted utilisation}}
     + \underbrace{\sum_{a} h_a \sum_{d} \phi_{a,d}}_{\text{holding fees}}
     \label{eq:cost}\\[4pt]
f_2 &= \max_{s} \sum_{d} d \cdot x_{s,d}
     - \min_{s} \sum_{d} d \cdot x_{s,d} + 1
     \label{eq:makespan}
\end{align}
where $c_{l,l'}$ is the composite transfer cost between locations (intra-town / inter-town / inter-city by geographic tier), $\sigma_{l,l',d}$ indicates a location change on day $d$, and $(\cdot)^{+} = \max(0,\cdot)$.

\paragraph{19 constraint families (C1--C19).}
\textbf{C1--C4}~\emph{Assignment}: each scene scheduled exactly once, at exactly one location, respecting precedence and time windows.
\textbf{C5--C7}~\emph{Capacity and availability}: daily shooting-hour cap; actor and location availability calendars.
\textbf{C8--C9}~\emph{Variable linking}: $v_{a,s,d}$ and $z_{l,s,d}$ couple scene scheduling to actor/location use; $w_{a,d}$ gives a scene-agnostic working-day indicator.
\textbf{C10--C13}~\emph{Geographic logistics}: location--town consistency; at most one active town per day (C11); no overnight town switches (C12); rolling consecutive-day labour limits (C13).
\textbf{C14}~\emph{Equipment sharing}: specialised equipment supply not exceeded per day.
\textbf{C15}~\emph{Booking and holding fees}: $\phi_{a,d}=1$ when actor is under contract but not shooting.
\textbf{C16--C17}~\emph{Scene coordination}: mandatory co-scheduling groups; minimum inter-scene separation.
\textbf{C18--C19}~\emph{Concurrency and rest}: location concurrency limits; mandatory rest windows for lead actors.

\subsection{Scene Interaction Graph (SIG)}
\label{sec:sig}

We build an undirected weighted graph on $\mathcal{S}$ with edge weights:
\begin{equation}
  w(s_i, s_j) = \alpha \cdot J_A(s_i, s_j)
              + \beta  \cdot J_L(s_i, s_j)
              + \gamma \cdot \frac{1}{1 + \bar{c}(s_i, s_j)}
  \label{eq:sig}
\end{equation}
where $J_A$, $J_L$ are Jaccard similarities of actor and location sets, $\bar{c}(s_i,s_j)$ is the average composite transfer cost, and $(\alpha,\beta,\gamma)=(0.5,0.3,0.2)$.
Scenes in the same mandatory co-scheduling group $G_k \in \mathcal{G}$ receive edge weight~1, ensuring they land in the same block.
Louvain community detection~\citep{blondel2008} on the SIG partitions $\mathcal{S}$ into shooting blocks $B_1,\ldots,B_K$, reducing the search space from $n!$ scene permutations to $K!$ block permutations.

\subsection{Bi-Level Matheuristic and Q-Learning AOS}
\label{sec:bilevel}

\textbf{Upper level.}
NSGA-II~\citep{deb2002} (population $N{=}20$, generations $G{=}40$) evolves permutations of block indices $\pi = (\pi_1, \ldots, \pi_K)$.
Non-dominated sorting and crowding distance maintain Pareto diversity.

\textbf{Lower level (global greedy reassembly).}
Given $\pi$, scenes are concatenated in block order; a greedy decoder assigns them to shooting days in that order, enforcing all C1--C19 constraints.
Days are shared across blocks, avoiding per-block cost overhead.
Each $\pi$ evaluates to an objective vector $(f_1, f_2)$.

\textbf{Q-learning AOS.}
State: 81 discretised states from four features (generation progress, population diversity, HV improvement rate, quality gap).
Action: one of 9 (crossover, mutation) operator pairs: $\{$PMX, OX, LAX$\} \times \{$swap, block-swap, ACM$\}$.
Reward: $\Delta\mathrm{HV}$ after applying the action.
The Q-table is updated online with learning rate $\eta{=}0.1$ and discount $\gamma_Q{=}0.9$.

\begin{algorithm}[t]
\caption{BDM-AOS: Bi-Level Decomposition Matheuristic with Adaptive Operator Selection}
\label{alg:bdmaos}
\begin{algorithmic}[1]
\REQUIRE MSSP instance; parameters $N$ (population), $G$ (generations), $p_c$, $p_m$
\ENSURE Approximated Pareto front of schedules
\STATE Build SIG via Eq.~\eqref{eq:sig}; $\{B_1,\ldots,B_K\} \leftarrow \mathrm{Louvain(SIG)}$
\STATE Initialise population $P = \{\pi^{(1)},\ldots,\pi^{(N)}\}$ with random block permutations
\STATE Evaluate each $\pi^{(i)}$ via greedy reassembly enforcing C1--C19
\FOR{$g = 1,\ldots,G$}
  \STATE $(cx, mut) \leftarrow \mathrm{AOS.select}(\mathrm{state})$ \hfill\COMMENT{Q-learning}
  \STATE $Q \leftarrow cx\text{-crossover}(P) \cup mut\text{-mutation}(P)$; evaluate $Q$
  \STATE $P \leftarrow \mathrm{NSGA\text{-}II\text{-}select}(P \cup Q)$
  \STATE $\mathrm{AOS.update}(\text{reward} = \Delta\mathrm{HV})$
\ENDFOR
\RETURN Non-dominated solutions from $P$
\end{algorithmic}
\end{algorithm}
